\thispagestyle{empty}
\vspace*{-2.6cm}
\begin{center}
  \begin{minipage}[c]{0.73\linewidth}
    \raggedright\includegraphics[height=1.5cm]{images_pdf/logo_Aix_Marseille_Universite.pdf}\quad\includegraphics[height=1.5cm]{images_pdf/logo_CNRS_fr_quadri_svg.pdf}\quad\includegraphics[height=1.5cm]{images_pdf/Logo_centrale_marseille.pdf}\quad\includegraphics[height=1.5cm]{images_pdf/logo_ED353_light_ligthblue_notext_4.pdf}
  \end{minipage}\hfill
  \begin{minipage}[c]{0.27\linewidth}
    \raggedleft\includegraphics[height=2.3cm]{images_pdf/logo_lma_picto_web_transp.pdf}
  \end{minipage}\hfill
\end{center}
\begin{flushleft}
%  \vspace{0.2cm}
% \LARGE
AIX-MARSEILLE UNIVERSITÉ\\
% \LARGE
CNRS, CENTRALE MARSEILLE\\
%  \vspace{0.2cm}
%%%%%%%%%%% \Large ED 353 SCIENCES POUR L'INGÉNIEUR: MÉCANIQUE, PHYSIQUE, MICRO ET NANOÉLECTRONIQUE\\
% \Large
  \small
ED 353 SCIENCES POUR L'INGÉNIEUR\\
MÉCANIQUE, PHYSIQUE, MICRO ET NANOÉLECTRONIQUE\\
%%%%%%%%%%%%%%  \vspace{0.2cm}
% \normalsize
  \vspace{0.5cm}
  LABORATOIRE DE MÉCANIQUE ET D'ACOUSTIQUE / UMR 7031 AMU-CNRS-ECM\\
  \vspace{0.5cm}
\end{flushleft}
%
\begin{center}
{
%\LARGE
\Large
\bf
Thèse de Doctorat\\
Spécialité :  Acoustique\\
\vspace{0.5cm}
présentée par Masaru NAGASO \ \ \ \includegraphics[height=0.7cm]{images_pdf/Masaru_Nagaso_name_in_Japanese2.pdf}
}\\[1.4cm]

{
%\LARGE
\Large
\bf
Étude de la propagation des ultrasons\\
dans un milieu fluide hétérogène en vue de la surveillance\\
en fonctionnement d'un réacteur nucléaire à caloporteur sodium

\vspace{1mm}

 ------------

\vspace{3mm}

\textsl{\textbf{Study of ultrasound wave propagation\\
in a heterogeneous fluid medium for the monitoring\\
\vspace{1.8mm}
of an operating sodium-based nuclear reactor}}
}

\end{center}

\vspace{7mm}

Devant être soutenue le 22 mai 2018 à 14h00 devant le jury composé de :\\
%
\begin{table}[hbt]
\makebox[1 \textwidth][c]{
\resizebox{1.0 \textwidth}{!}{

  \centering
  \footnotesize
   \begin{tabularx}{1.0\textwidth}{lll}
   \toprule
    Mounsif Ech-Cherif El-Kettani & Professeur des Universités, Université du Havre            & Rapporteur  \\
    Michel Darmon                 & Ingénieur Habilité à Diriger des Recherches, CEA Saclay    & Rapporteur  \\
    Stéphane Lanteri              & Directeur de Recherche INRIA, Nice Sophia Antipolis        & Examinateur \\
    Antoine Gerschenfeld          & Ingénieur CEA Saclay, Docteur ès Sciences                  & Examinateur \\
    Dominique Habault             & Directrice de Recherche CNRS, LMA Marseille                & Examinatrice \\
    Dimitri Komatitsch            & Directeur de Recherche CNRS, LMA Marseille                 & Directeur de thèse \\
    Joseph Moysan                 & Professeur des Universités, Université d'Aix-Marseille, LMA  & Co-directeur de thèse \\
    Christian Lhuillier           & Ingénieur CEA Cadarache                                    & Encadrant CEA, invité \\ \bottomrule
   \end{tabularx}
}
}
\end{table}
\vfill
%\normalsize\rm
\small\rm
\begin{center}
Masaru NAGASO\\
Laboratoire de Mécanique et d'Acoustique UMR 7031 AMU-CNRS-ECM\\
4 impasse Nikola Tesla, 13453 Marseille cedex 13, France\\
Email: masaru.nagaso@cea.fr and mnsaru22@gmail.com
\end{center}
\newpage
\thispagestyle{empty}
\normalsize
